\chapter{字符串}

\section{字符串匹配}

\subsection{KMP}

在解题函数内定义下列 Lambda。调用 \texttt{KMP(s,t)}，\texttt{s} 为文本，\texttt{t} 为非空模式串。每个匹配起点单独输出一行，下标从 $1$ 开始。

\cppfile{code/string/kmp.cpp}

\subsection{Z 函数}

调用 \texttt{zFunction(s)}。对 $0\le i<|s|$，返回值 \texttt{z[i]} 是后缀 \texttt{s[i..]} 与整个字符串的最长公共前缀长度。

\cppfile{code/string/Z.cpp}

\subsection{字符串双哈希}

选用本册旧版 \texttt{MInt}，调用 \texttt{init(n)} 预处理到最大串长。用 \texttt{String h(s)} 构造哈希。\texttt{get()} 返回整串哈希，\texttt{substring(l,r)} 返回下标从 $0$ 开始的闭区间哈希。\texttt{modify(i,c)} 只返回替换后的哈希对，保留原对象。要与 \texttt{get()} 比较，分别减去对应的 \texttt{val1[n]} 和 \texttt{val2[n]}。

\cppfile{code/string/hash.cpp}

\section{字符串数据结构}

\subsection{Trie 字典树}

用 \texttt{Trie t(n)} 预留节点，包含根和所有可能插入的字符。字符串只用大小写字母和数字。用 \texttt{insert(s)} 插入，\texttt{query(s)} 返回整串出现次数。调用 \texttt{del(s)} 前确认该串存在。

\cppfile{code/string/trie.cpp}

\subsection{AC 自动机}

把所有非空小写模式串用 \texttt{insert(s)} 插入，调用一次 \texttt{build()}，再用 \texttt{query(t)} 统计至少出现一次的模式串数量。重复插入的模式按次数计。每次查询会消耗计数，多次独立查询需要重建。

\cppfile{code/string/ac.cpp}

\section{常考进阶}

\subsection{Manacher}

调用 \texttt{manacher(s)}，输入字符避开分隔符 \texttt{\#}。返回的 \texttt{r[i]} 对应插入分隔符后的中心下标，原串的最长回文子串长度取 \texttt{max(r)-1}。

\cppfile{code/string/manacher.cpp}
